\par\Needspace{6\baselineskip}\medskip
\hypertarget{native-132-article097}{}\pdfbookmark[2]{97.}{article-097}
\begin{center}97.\end{center}
Zur Erläuterung der seit Art. 88 entwickelten Methode greifen wir die beiden Beispiele aus Art. 87 wieder auf. Es braucht kaum erwähnt zu werden, dass die bisher verwendete Bedeutung des Hilfswinkels $\omega$ von der Bedeutung zu unterscheiden ist, die dasselbe Zeichen in Art. 86 und 87 hatte.

I. Im ersten Beispiel gilt $f=3^0\,47'\,26''{,}865$. Ferner ergibt sich
\[
\begin{gathered}
\log\dfrac{r'}r=9{,}991\,4599,\qquad
\log\operatorname{tang}(45^0+\omega)=9{,}997\,8649\,75,\\
\omega=-8'\,27''{,}006.
\end{gathered}
\]
\noindent Daraus folgt nach Art. 89

\begin{center}
\begin{tabular}{@{}ll@{}}
\makebox[17em][l]{$\log\sin\tfrac12 f^2$\hfill .\,.\,.\,.\,.\,.\,.\,.\,.\,.} & $7{,}038\,9972$\\[.3em]
\makebox[17em][l]{$\log\cos f$\hfill .\,.\,.\,.\,.\,.\,.\,.\,.\,.\,.\,.} & $9{,}999\,0488$\\[.3em]
\hline
 & $7{,}039\,9484$\\[.3em]
\multicolumn{2}{l}{$=\log0{,}001\,0963\,480$}\\[.3em]
\end{tabular}
\end{center}

\begin{center}
\begin{tabular}{@{}ll@{}}
\makebox[17em][l]{$\log\operatorname{tang}2\omega^2$\hfill .\,.\,.\,.\,.\,.\,.\,.} & $5{,}383\,2428$\\[.3em]
\makebox[17em][l]{$\log\cos f$\hfill .\,.\,.\,.\,.\,.\,.\,.\,.\,.\,.} & $9{,}999\,0488$\\[.3em]
\hline
 & $5{,}384\,1940$\\[.3em]
\multicolumn{2}{l}{$=\log0{,}000\,0242\,211$}\\[.3em]
\end{tabular}
\end{center}

\noindent\SourceLeaf{133}und folglich $l=0{,}001\,1205\,691$, $\tfrac56+l=0{,}834\,4539$. Ferner ergibt sich $\log kt=9{,}576\,6974$

\begin{center}
\begin{tabular}{@{}ll@{}}
\makebox[17em][l]{$2\log kt$\hfill .\,.\,.\,.\,.\,.\,.\,.\,.\,.\,.\,.\,.\,.} & $9{,}153\,3948$\\[.3em]
\makebox[17em][l]{$\mathrm{C.}\tfrac32\log rr'$\hfill .\,.\,.\,.\,.\,.\,.\,.\,.\,.\,.} & $9{,}020\,5181$\\[.3em]
\makebox[17em][l]{$\mathrm{C.}\log8\cos f^3$\hfill .\,.\,.\,.\,.\,.\,.\,.\,.} & $9{,}099\,7636$\\[.3em]
\hline
\makebox[17em][l]{$\log mm$\hfill .\,.\,.\,.\,.\,.\,.\,.\,.\,.\,.\,.\,.\,.} & $7{,}273\,6765$\\[.3em]
\makebox[17em][l]{$\log(\tfrac56+l)$\hfill .\,.\,.\,.\,.\,.\,.\,.\,.\,.\,.\,.} & $9{,}921\,4023$\\[.3em]
\hline
 & $7{,}352\,2742$\\[.3em]
\end{tabular}
\end{center}

\noindent Der Näherungswert ist also $h=0{,}002\,2504\,7$. Ihm entspricht in unserer Tafel II der Wert $\log yy=0{,}002\,1633$. Somit ergibt sich $\log\dfrac{mm}{yy}=7{,}271\,5132$ oder $\dfrac{mm}{yy}=0{,}001\,8685\,87$, woraus nach Formel 16 $x=0{,}000\,7480\,179$ folgt. Da $\xi$ nach Tafel III gänzlich vernachlässigbar ist, brauchen die gefundenen Werte für $h$, $y$ und $x$ keine Korrektur. Die Elemente werden nun wie folgt bestimmt.

\begin{center}
\begin{tabular}{@{}ll@{}}
\makebox[17em][l]{$\log x$\hfill .\,.\,.\,.\,.\,.\,.\,.\,.\,.\,.\,.\,.\,.\,.\,.} & $6{,}873\,9120$\\[.3em]
\makebox[17em][l]{$\log\sin\tfrac12 g$\hfill .\,.\,.\,.\,.\,.\,.\,.\,.\,.\,.\,.} & $8{,}436\,9560$\\[.3em]
\end{tabular}
\end{center}

\[
\begin{gathered}
\tfrac12 g=1^0\,34'\,2''{,}0286,\qquad
\tfrac12(f+g)=3^0\,27'\,45''{,}4611,\\
\tfrac12(f-g)=19'\,41''{,}4039.
\end{gathered}
\]
\Needspace{19\baselineskip}
\noindent Nach den Formeln 27, 28, 29 und 30 erhält man daher

\begin{center}
\begin{tabular}{@{}ll@{}}
\makebox[17em][l]{$\log\operatorname{tang}2\omega$\hfill .\,.\,.\,.\,.\,.\,.\,.\,.\,.\,.} & $7{,}691\,6214_n$\\[.3em]
\makebox[17em][l]{$\log\cos\tfrac12(f+g)$\hfill .\,.\,.\,.\,.\,.\,.\,.} & $9{,}999\,2065$\\[.3em]
\makebox[17em][l]{$\log\cos\tfrac12(f-g)$\hfill .\,.\,.\,.\,.\,.\,.\,.} & $9{,}999\,9929$\\[.3em]
\hline\noalign{\vskip .3em}
\makebox[17em][l]{$\log P\sin\tfrac12(F-G)$\hfill .\,.\,.\,.\,.} & $7{,}690\,8279_n$\\[.3em]
\makebox[17em][l]{$\log P\cos\tfrac12(F-G)$\hfill .\,.\,.\,.\,.} & $8{,}781\,0240$\\[.3em]
\hline\noalign{\vskip .3em}
$\tfrac12(F-G)=$ & $-4^0\,38'\,41''{,}54$\\[.3em]
$\tfrac12(F+G)=$ & $319\quad21\quad38{,}05$\\[.3em]
\hline
$F=$ & $314\quad42\quad56{,}51$\\[.3em]
$v=$ & $310\quad55\quad29{,}64$\\[.3em]
$v'=$ & $318\quad30\quad23{,}37$\\[.3em]
$G=$ & $324\quad0\quad19{,}59$\\[.3em]
$E=$ & $320\quad52\quad15{,}53$\\[.3em]
$E'=$ & $327\quad8\quad23{,}65$\\[.3em]
\end{tabular}
\end{center}

\begin{center}
\begin{tabular}{@{}ll@{}}
\makebox[17em][l]{$\mathrm{C.}\log\cos2\omega$\hfill .\,.\,.\,.\,.\,.\,.\,.\,.\,.} & $0{,}000\,0052$\\[.3em]
\makebox[17em][l]{$\log\sin\tfrac12(f+g)$\hfill .\,.\,.\,.\,.\,.\,.\,.} & $8{,}781\,0188$\\[.3em]
\makebox[17em][l]{$\log\sin\tfrac12(f-g)$\hfill .\,.\,.\,.\,.\,.\,.\,.} & $7{,}757\,9709$\\[.3em]
\hline\noalign{\vskip .3em}
\makebox[17em][l]{$\log Q\sin\tfrac12(F+G)$\hfill .\,.\,.\,.\,.} & $7{,}691\,6143_n$\\[.3em]
\makebox[17em][l]{$\log Q\cos\tfrac12(F+G)$\hfill .\,.\,.\,.\,.} & $7{,}757\,9761$\\[.3em]
\hline\noalign{\vskip .3em}
\makebox[17em][l]{$\log P=\log R\cos\tfrac12\varphi$\hfill .\,.\,.} & $8{,}782\,4527$\\[.3em]
\makebox[17em][l]{$\log Q=\log R\sin\tfrac12\varphi$\hfill .\,.\,.} & $7{,}877\,8355$\\[.3em]
\hline\noalign{\vskip .3em}
$\text{Daraus }\tfrac12\varphi=$ & $7^0\,6'\,0''{,}935$\\[.3em]
$\varphi=$ & $14\quad12\quad1{,}87$\\[.3em]
\makebox[17em][l]{$\log R$\hfill .\,.\,.\,.\,.\,.\,.\,.\,.\,.\,.\,.\,.\,.\,.} & $8{,}785\,7960$\\[.3em]
\end{tabular}
\end{center}

\begin{center}\small Zur Kontrolle der Rechnung\end{center}

\begin{center}
\begin{tabular}{@{}ll@{}}
\makebox[17em][l]{$\tfrac12\log2\cos f$\hfill .\,.\,.\,.\,.\,.\,.\,.\,.\,.\,.} & $0{,}150\,0394$\\[.3em]
\makebox[17em][l]{$\tfrac12\log(l+x)=\log\dfrac m y$\hfill .\,.\,.\,.} & $8{,}635\,7566$\\[.3em]
\noalign{\vskip .3em}\hline
 & $8{,}785\,7960$\\[.3em]
\end{tabular}
\end{center}

\SourceLeaf{134}

\begin{center}
\begin{tabular}{@{}ll@{}}
\makebox[17em][l]{$\tfrac12\log rr'$\hfill .\,.\,.\,.\,.\,.\,.\,.\,.\,.\,.\,.\,.\,.} & $0{,}326\,4939$\\[.3em]
\makebox[17em][l]{$\log\sin f$\hfill .\,.\,.\,.\,.\,.\,.\,.\,.\,.\,.\,.\,.\,.\,.} & $8{,}820\,2909$\\[.3em]
\makebox[17em][l]{$\mathrm{C.}\log\sin g$\hfill .\,.\,.\,.\,.\,.\,.\,.\,.\,.\,.\,.} & $1{,}262\,1765$\\[.3em]
\hline
\makebox[17em][l]{$\log b$\hfill .\,.\,.\,.\,.\,.\,.\,.\,.\,.\,.\,.\,.\,.\,.\,.\,.} & $0{,}408\,9613$\\[.3em]
\makebox[17em][l]{$\log\cos\varphi$\hfill .\,.\,.\,.\,.\,.\,.\,.\,.\,.\,.\,.\,.\,.} & $9{,}986\,5224$\\[.3em]
\hline
\makebox[17em][l]{$\log p$\hfill .\,.\,.\,.\,.\,.\,.\,.\,.\,.\,.\,.\,.\,.\,.\,.\,.} & $0{,}395\,4837$\\[.3em]
\makebox[17em][l]{$\log a$\hfill .\,.\,.\,.\,.\,.\,.\,.\,.\,.\,.\,.\,.\,.\,.\,.\,.} & $0{,}422\,4389$\\[.3em]
\end{tabular}
\end{center}

\begin{center}
\begin{tabular}{@{}ll@{}}
\makebox[17em][l]{$\log k$\hfill .\,.\,.\,.\,.\,.\,.\,.\,.\,.\,.\,.\,.\,.\,.\,.\,.} & $3{,}550\,0066$\\[.3em]
\makebox[17em][l]{$\tfrac32\log a$\hfill .\,.\,.\,.\,.\,.\,.\,.\,.\,.\,.\,.\,.\,.\,.\,.} & $0{,}633\,6584$\\[.3em]
\hline
 & $2{,}916\,3482$\\[.3em]
\makebox[17em][l]{$\log t$\hfill .\,.\,.\,.\,.\,.\,.\,.\,.\,.\,.\,.\,.\,.\,.\,.\,.} & $1{,}341\,1160$\\[.3em]
\hline
 & $4{,}257\,4642$\\[.3em]
\end{tabular}
\end{center}

\noindent Die mittlere tägliche Bewegung ist also $=824''{,}7990$. Die mittlere Bewegung in der Zeit $t=18091''{,}07=5^0\,1'\,31''{,}07$.

\begin{center}
\begin{tabular}{@{}ll@{}}
\makebox[17em][l]{$\log\sin\varphi$\hfill .\,.\,.\,.\,.\,.\,.\,.\,.\,.\,.\,.\,.\,.} & $9{,}389\,7262$\\[.3em]
\makebox[17em][l]{$\log206265$\hfill .\,.\,.\,.\,.\,.\,.\,.\,.\,.\,.\,.} & $5{,}314\,4251$\\[.3em]
\hline
\makebox[17em][l]{$\log e\text{ in Bogensekunden}$\hfill .\,.\,.\,.\,.\,.\,.\,.} & $4{,}704\,1513$\\[.3em]
\makebox[17em][l]{$\mathop{\mathrm{lo}\bar{\mathrm{g}}}\nolimits\sin E$\hfill .\,.\,.\,.\,.\,.\,.\,.\,.\,.\,.\,.\,.\,.} & $9{,}800\,0767_n$\\[.3em]
\makebox[17em][l]{$\log\sin E'$\hfill .\,.\,.\,.\,.\,.\,.\,.\,.\,.\,.\,.\,.} & $9{,}734\,4714_n$\\[.3em]
\hline
\makebox[17em][l]{$\log e\sin E$\hfill .\,.\,.\,.\,.\,.\,.\,.\,.\,.\,.\,.\,.} & $4{,}504\,2280_n$\\[.3em]
\makebox[17em][l]{$\log e\sin E'$\hfill .\,.\,.\,.\,.\,.\,.\,.\,.\,.\,.\,.} & $4{,}438\,6227_n$\\[.3em]
\end{tabular}
\end{center}

\[
\begin{aligned}
e\sin E&=-31932''{,}14=-8^0\,52'\,12''{,}14\\
e\sin E'&=-27455{,}08=-7\quad37\quad35{,}08
\end{aligned}
\]
\begin{center}Daraus die mittlere Anomalie\end{center}

\begin{center}
\begin{tabular}{@{}ll@{}}
$\text{für den ersten Ort}=$ & $329^0\,44'\,27''{,}67$\\[.3em]
$\text{für den zweiten}=$ & $334\quad45\quad58{,}73$\\[.3em]
\hline
$\text{Differenz}=$ & $5\quad1\quad31{,}06$\\[.3em]
\end{tabular}
\end{center}

II. Im zweiten Beispiel erhält man $f=31^0\,27'\,38''{,}32$, $\omega=-21'\,50''{,}565$, $l=0{,}086\,3565\,9$, $\log mm=9{,}353\,0651$, und der Quotient $\dfrac{mm}{\frac56+l}$ liefert den Näherungswert $h=0{,}245\,1454$. Dem entspricht in Tafel II der Wert $\log yy=0{,}172\,2663$, aus dem sich $\dfrac{mm}{yy}=0{,}151\,6347\,7$ und $x=0{,}065\,2781\,8$ ergeben. Damit entnimmt man Tafel III den Wert $\xi=0{,}000\,2531$. Verwendet man diesen Wert, so erhält man die korrigierten Werte $h=0{,}245\,0779$, $\log yy=0{,}172\,2303$, $\dfrac{mm}{yy}=0{,}151\,6473\,5$, $x=0{,}065\,2907\,6$ und $\xi=0{,}000\,2532$. Wiederholte man die Rechnung mit diesem Wert von $\xi$, der sich vom vorigen nur um eine Einheit an der siebten Dezimalstelle unterscheidet, so würden sich $h$, $\log yy$ und $x$ nicht merklich ändern. Der gefundene Wert von $x$ ist daher bereits der wahre Wert, und man kann unmittelbar zur Bestimmung der Elemente übergehen. Darauf gehen wir hier nicht näher ein, da sie genau wie im vorigen Beispiel verläuft.

III. Auch den anderen Fall, in dem $\cos f$ negativ ist, wollen wir an einem Beispiel erläutern. Es sei $v'-v=224^0\,0'\,0''$ oder $f=112^0\,0'\,0''$, ferner $\log r=0{,}139\,4892$, $\log r'=0{,}397\,8794$ und $t=206{,}80919$ Tage. Hier erhält man $\omega=+4^0\,14'\,43''{,}82$, $L=1{,}894\,2298$, $\log MM=0{,}672\,4333$ und als ersten Näherungswert $\log H=0{,}646\,7603$. Durch Lösen der Gleichung $15^*$ erhält man $Y=1{,}591\,432$ und daraus $x=0{,}037\,037$, wozu in Tafel III der Wert $\xi=0{,}000\,0801$ gehört. \SourceLeaf{135}Daraus ergeben sich die korrigierten Werte $\log H=0{,}646\,7931$, $Y=1{,}591\,5107$, $x=0{,}037\,2195$ und $\xi=0{,}000\,0809$. Wiederholt man die Rechnung mit diesem Wert von $\xi$, so erhält man $x=0{,}037\,2213$. Dieser Wert bedarf keiner weiteren Korrektur, da sich daraus derselbe Wert von $\xi$ ergibt. Anschließend erhält man $\tfrac12 g=11^0\,7'\,25''{,}40$ und daraus, ebenso wie in Beispiel I,

\begin{center}
\begin{tabular}{@{}ll@{}}
$\tfrac12(F-G)=$ & $3^0\,33'\,53''{,}62$\\[.3em]
$\tfrac12(F+G)=$ & $8\quad26\quad6{,}38$\\[.3em]
$F=$ & $12\quad0\quad0{,}00$\\[.3em]
$v=$ & $-100\quad0\quad0{,}00$\\[.3em]
$v'=$ & $+124\quad0\quad0{,}00$\\[.3em]
$G=$ & $4\quad52\quad12{,}76$\\[.3em]
$E=$ & $-17\quad22\quad38{,}04$\\[.3em]
$E'=$ & $+27\quad7\quad3{,}56$\\[.3em]
\end{tabular}
\end{center}

\begin{center}
\begin{tabular}{@{}ll@{}}
\makebox[17em][l]{$\log P=\log R\cos\tfrac12\varphi$\hfill .\,.\,.\,.} & $9{,}970\,0508$\\[.3em]
\makebox[17em][l]{$\log Q=\log R\sin\tfrac12\varphi$\hfill .\,.\,.\,.} & $9{,}858\,0553$\\[.3em]
$\tfrac12\varphi=$ & $37^0\,41'\,34''{,}27$\\[.3em]
$\varphi=$ & $75\quad23\quad8{,}54$\\[.3em]
\makebox[17em][l]{$\log R$\hfill .\,.\,.\,.\,.\,.\,.\,.\,.\,.\,.\,.\,.\,.\,.\,.\,.} & $0{,}071\,7097$\\[.3em]
\end{tabular}
\end{center}

\Needspace{6\baselineskip}
\begin{center}\small Zur Kontrolle der Rechnung erhält man\end{center}

\begin{center}
\begin{tabular}{@{}ll@{}}
\makebox[17em][l]{$\log\dfrac M Y\BareRoot{(-2\cos f)}$\hfill .\,.\,.\,.\,.\,.\,.} & $0{,}071\,7097$\\[.3em]
\end{tabular}
\end{center}

\noindent Bei Bahnen von so großer Exzentrizität lässt sich der Winkel $\varphi$ nach Formel $19^*$ etwas genauer berechnen. Sie liefert in unserem Beispiel $\varphi=75^0\,23'\,8''{,}57$. Auch die Exzentrizität $e$ lässt sich nach der Formel $1-2\sin(45^0-\tfrac12\varphi)^2$ genauer bestimmen als nach $\sin\varphi$. Die erste Formel ergibt $e=0{,}967\,6463\,0$.

Ferner ergibt Formel 1 den Wert $\log b=0{,}657\,6611$ und daraus $\log p=0{,}0595967$, $\log a=1{,}255\,7255$ sowie für den Logarithmus des Abstands im Perihel
\[
\begin{aligned}
{}={}&\log\dfrac p{1+e}=\log(a(1-e))\\
{}={}&\log(b\operatorname{tang}(45^0-\tfrac12\varphi))=9{,}765\,6496.
\end{aligned}
\]

Bei Bahnen, die der Parabel so nahekommen, gibt man gewöhnlich den Zeitpunkt des Durchgangs durch das Perihel anstelle der Epoche der mittleren Anomalie an. Sind die Elemente bekannt, so lassen sich die Zeitabstände zwischen diesem Zeitpunkt und den Zeitpunkten der beiden gegebenen Orte nach der in Art. 41 angegebenen Methode bestimmen. Ihre Differenz oder Summe, je nachdem, ob das Perihel außerhalb oder zwischen den beiden gegebenen Orten liegt, dient zur Kontrolle der Rechnung, da sie mit der Zeit $t$ übereinstimmen muss. — Die Zahlen dieses dritten Beispiels beruhten auf den Elementen, die im Beispiel aus Art. 38 und 43 vorausgesetzt wurden. Dieses Beispiel lieferte sogar unseren ersten Ort. Die geringfügigen Abweichungen der hier ermittelten Elemente rührten allein von der begrenzten Genauigkeit der logarithmischen und trigonometrischen Tafeln her.
